\chapter{CONCLUSÃO}

O desenvolvimento do CreditMaster evidenciou que a utilização de ferramentas de Inteligência Artificial Generativa arquitetadas e executadas de forma inteiramente local (\textit{on-premise}) constitui uma resposta robusta, escalável e segura aos desafios de modernização do setor bancário. A arquitetura híbrida delineada ao longo do trabalho resolve, com eficácia, dois problemas críticos da adoção em massa de LLMs em instituições financeiras: a manutenção incondicional do sigilo dos dados financeiros dos clientes e a eliminação sistêmica do risco de \textit{alucinações}.

Ao conjugar o \textit{Fine-Tuning} de modelos avançados — como o Llama 3.1 8B otimizado através de técnicas QLoRA e empacotamento GGUF — com um sistema maduro de \textit{Retrieval-Augmented Generation} (RAG), a solução alcançou resultados satisfatórios. O modelo não apenas mimetizou com perfeição o tom professoral e a terminologia estrita dos comitês de risco, como passou a basear cada decisão gerada nas orientações mais recentes de manuais e políticas internas da instituição. O CreditMaster, portanto, posiciona-se como uma evolução plausível da esteira clássica de crédito, transferindo o trabalho braçal da escrita e validação humana inicial para um \textit{microsserviço headless} ultrarrápido operado sobre recursos de \textit{hardware} acessíveis.

\textbf{[ANOTAÇÃO: Preciso complementar isso aqui dizendo se consegui bater as metas de precisão e qual foi a minha taxa de sucesso nos pareceres durante os testes.]}

\section{Trabalhos Futuros}
\textbf{[ANOTAÇÃO: Listar algumas ideias pro futuro: 
1. Tentar rodar modelos maiores (tipo 70B quantizado) se arrumar hardware.
2. Fazer uma integração de verdade com sistemas do banco via API.
3. Colocar as resoluções do Bacen no banco vetorial também.
4. Fazer um front-end bonitão pra alguém revisar os pareceres antes de publicar.]}
